% Options for packages loaded elsewhere
\PassOptionsToPackage{unicode}{hyperref}
\PassOptionsToPackage{hyphens}{url}
\documentclass[
  11pt,
  a4paper,
  fontset=fandol]{ctexart}
\usepackage{xcolor}
\usepackage[margin=2cm]{geometry}
\usepackage{amsmath,amssymb}
\setcounter{secnumdepth}{-\maxdimen} % remove section numbering
\usepackage{iftex}
\ifPDFTeX
  \usepackage[T1]{fontenc}
  \usepackage[utf8]{inputenc}
  \usepackage{textcomp} % provide euro and other symbols
\else % if luatex or xetex
  \usepackage{unicode-math} % this also loads fontspec
  \defaultfontfeatures{Scale=MatchLowercase}
  \defaultfontfeatures[\rmfamily]{Ligatures=TeX,Scale=1}
\fi
\usepackage{lmodern}
\ifPDFTeX\else
  % xetex/luatex font selection
\fi
% Use upquote if available, for straight quotes in verbatim environments
\IfFileExists{upquote.sty}{\usepackage{upquote}}{}
\IfFileExists{microtype.sty}{% use microtype if available
  \usepackage[]{microtype}
  \UseMicrotypeSet[protrusion]{basicmath} % disable protrusion for tt fonts
}{}
\usepackage{setspace}
\makeatletter
\@ifundefined{KOMAClassName}{% if non-KOMA class
  \IfFileExists{parskip.sty}{%
    \usepackage{parskip}
  }{% else
    \setlength{\parindent}{0pt}
    \setlength{\parskip}{6pt plus 2pt minus 1pt}}
}{% if KOMA class
  \KOMAoptions{parskip=half}}
\makeatother
\usepackage{longtable,booktabs,array}
\usepackage{caption}
\captionsetup[table]{skip=6pt}
\newcounter{none} % for unnumbered tables
\usepackage{calc} % for calculating minipage widths
% Correct order of tables after \paragraph or \subparagraph
\usepackage{etoolbox}
\makeatletter
\patchcmd\longtable{\par}{\if@noskipsec\mbox{}\fi\par}{}{}
\makeatother
\setlength{\emergencystretch}{3em} % prevent overfull lines
\providecommand{\tightlist}{%
  \setlength{\itemsep}{0pt}\setlength{\parskip}{0pt}}
\usepackage{amsmath,amssymb}
\usepackage{fancyhdr}
\usepackage{needspace}
\setlength{\headheight}{15pt}
\pagestyle{fancy}
\fancyhf{}
\fancyhead[L]{\small 高等代数期中计算复习}
\fancyhead[R]{\small 行列式\quad 秩\quad 方程组}
\fancyfoot[C]{\thepage}
\renewcommand{\headrulewidth}{0.3pt}
\setlength{\parskip}{0.25em}
\raggedbottom
\clubpenalty=10000
\widowpenalty=10000
\displaywidowpenalty=10000
\makeatletter
\renewcommand{\maketitle}{%
  \thispagestyle{plain}%
  \begin{center}
    {\Large\bfseries\@title\par}\vspace{0.4em}
    {\small\@date\par}
  \end{center}\vspace{0.4em}
}
\makeatother
\AtBeginDocument{\hypersetup{colorlinks=true,linkcolor=black,urlcolor=black}}
\usepackage{bookmark}
\IfFileExists{xurl.sty}{\usepackage{xurl}}{} % add URL line breaks if available
\urlstyle{same}
% fallback for those not using the hyperref driver hyperxmp:
\makeatletter
\@ifundefined{xmpquote}{\newcommand{\xmpquote}[1]{#1}}{}
\makeatother
\hypersetup{
  pdftitle={高等代数期中计算复习},
  hidelinks,
  pdfcreator={LaTeX via pandoc}}

\title{高等代数期中计算复习}
\author{}
\date{2026年11月10日}

\begin{document}
\maketitle

\setstretch{1.08}
同一个矩阵，既能用来求秩，也能用来找极大无关组和求方程组的解。我们先把这一组计算连起来，再处理含参数方程组和行列式。

前五节共有六个短检查，每次先自己写出结果，最后到第六节核对。卡住三分钟就做记号，先读下一步。第一遍的目标是把整条计算过程走通，最后再重做有记号的地方。

下面是约两小时的阅读安排。每完成一节，勾掉一项；第七节是选读补充。

\begin{itemize}
\tightlist
\item[$\square$]
  第一节 行变换 10分钟
\item[$\square$]
  第二节 秩定理与向量组 25分钟
\item[$\square$]
  第三节 方程组 25分钟
\item[$\square$]
  第四节 参数分类 20分钟
\item[$\square$]
  第五节 行列式 25分钟
\item[$\square$]
  第六节 核对答案与回看 15分钟
\end{itemize}

\Needspace{7\baselineskip}

\section{1
先完成一次行消元}\label{1-ux5148ux5b8cux6210ux4e00ux6b21ux884cux6d88ux5143}

消元用到三种初等行变换。先分清它们怎样影响行列式、秩和解集，再开始计算。

{\def\LTcaptype{none} % do not increment counter
\begin{longtable}[]{@{}llll@{}}
\toprule\noalign{}
操作 & 行列式的变化（方阵时） & 矩阵的秩 & 对增广矩阵操作时的解集 \\
\midrule\noalign{}
\endhead
\bottomrule\noalign{}
\endlastfoot
交换两行 & 变号 & 不变 & 不变 \\
一行乘非零数 k & 乘 k & 不变 & 不变 \\
一行加上另一行的 k 倍 & 不变 & 不变 & 不变 \\
\end{longtable}
}

这里的``另一行''不能是被替换的行。若有参数，作除法以前，先检查分母是否可能为零。对方程组消元时，常数列必须跟着整行一起变换。

前三节都使用这个矩阵 A：

\[A=\begin{pmatrix}
1&2&-1&1\\
2&4&1&-1\\
1&2&2&-2
\end{pmatrix}.\]

第一行首项是
1，先用它消去下面两行的第一项：第二行减第一行的两倍，第三行减第一行。

\[\begin{pmatrix}
1&2&-1&1\\2&4&1&-1\\1&2&2&-2
\end{pmatrix}
\xrightarrow[R_3\leftarrow R_3-R_1]{R_2\leftarrow R_2-2R_1}
\begin{pmatrix}
1&2&-1&1\\0&0&3&-3\\0&0&3&-3
\end{pmatrix}.\]

第二、三行现在相同，第三行减去第二行：

\[\begin{pmatrix}
1&2&-1&1\\0&0&3&-3\\0&0&3&-3
\end{pmatrix}
\xrightarrow{R_3\leftarrow R_3-R_2}
C=\begin{pmatrix}
1&2&-1&1\\0&0&3&-3\\0&0&0&0
\end{pmatrix}.\]

这已经是行阶梯形。各非零行的第一个非零元素叫主元，主元所在的列叫主元列。\textbf{只求秩时，化到这里就够了，不必继续把每个主元都化成
1。}

\textbf{短检查1}
这个阶梯形矩阵有几个非零行？主元分别在哪几列？先写下这两个结果。

\Needspace{7\baselineskip}

\section{2
用秩定理读出向量组的信息}\label{2-ux7528ux79e9ux5b9aux7406ux8bfbux51faux5411ux91cfux7ec4ux7684ux4fe1ux606f}

上一步得到两个非零行。秩定理把这个结果与行向量、列向量联系起来：

\[\boxed{\operatorname{rank}A=\text{行秩}(A)=\text{列秩}(A).}\]

行秩是行向量组的秩，列秩是列向量组的秩；向量组的秩，就是极大线性无关组所含向量的个数。行阶梯形的非零行线性无关，因此本例的秩是
2。A 经过初等行变换得到 C，所以 A 的秩也是 2。

\Needspace{7\baselineskip}

\subsection{从主元列选出原向量}\label{ux4eceux4e3bux5143ux5217ux9009ux51faux539fux5411ux91cf}

要选原列向量组的极大无关组，先记主元列号，再回到原矩阵取列。本例主元在第1、3列，因此可选

\[\boldsymbol\alpha_1=\begin{pmatrix}1\\2\\1\end{pmatrix},
\qquad
\boldsymbol\alpha_3=\begin{pmatrix}-1\\1\\2\end{pmatrix}.\]

这两个向量来自原矩阵 A。行变换以后，同一列的分量通常变了，不能把 C
的第1、3列直接当作原向量组中的向量交上去。

为什么这两列足够？先看原矩阵中的一个二阶子式：

\[\begin{vmatrix}1&-1\\2&1\end{vmatrix}=3\ne0.\]

它说明第1、3列线性无关，也说明矩阵的秩至少为
2。再结合已经算出的阶梯形，知道秩至多为 2，因此秩恰好是
2。这两列便构成一个极大线性无关组。

\textbf{找到一个非零二阶子式，只能得到秩至少为 2。} 要断定秩就是
2，还需要上界。本例也可以直接看到原第三行等于原第二行减原第一行，所以三行相关，秩不超过
2。

\Needspace{7\baselineskip}

\subsection{从列关系写出线性表示}\label{ux4eceux5217ux5173ux7cfbux5199ux51faux7ebfux6027ux8868ux793a}

原列之间还能怎样表示？继续看刚才的阶梯形：第2列是第1列的两倍，第4列是第3列的相反数。可逆的行变换保持列之间具有同一组系数的线性关系，因此原矩阵中也有

\[\boldsymbol\alpha_2=2\boldsymbol\alpha_1,
\qquad
\boldsymbol\alpha_4=-\boldsymbol\alpha_3.\]

所以这四个列向量都能由第1、3列表出。极大无关组可以不唯一，但每个极大无关组都恰好含有两个向量。

\textbf{短检查2}
原矩阵的第2、4列，能否组成另一个极大线性无关组？说明你使用的两条列关系。

``两个向量组秩相同''不能推出它们等价。比如二维向量 (1,0) 与 (0,1)
各自组成的单向量组，秩都是
1，却不能互相线性表示。\textbf{等价要求双方能互相表出。}

\Needspace{7\baselineskip}

\subsection{用子式判断含参数矩阵的秩}\label{ux7528ux5b50ux5f0fux5224ux65adux542bux53c2ux6570ux77e9ux9635ux7684ux79e9}

非零矩阵的秩，还等于它的非零子式的最高阶数；零矩阵的秩规定为 0。对 n
阶方阵，先算行列式：非零时秩就是
n；为零时，代入对应参数值，再消元或检查低阶子式。

\[\det A\ne0\Longleftrightarrow\operatorname{rank}A=n
\quad\text{仅在 }A\text{ 是 }n\text{ 阶方阵时使用}.\]

\Needspace{7\baselineskip}

\section{3
从同一个矩阵求出全部解}\label{3-ux4eceux540cux4e00ux4e2aux77e9ux9635ux6c42ux51faux5168ux90e8ux89e3}

现在把 A
的列当作系数列，刚才的消元就能用来求解。先看齐次方程组，再只改变常数列。

\Needspace{7\baselineskip}

\subsection{齐次方程组与基础解系}\label{ux9f50ux6b21ux65b9ux7a0bux7ec4ux4e0eux57faux7840ux89e3ux7cfb}

齐次方程组的右端全为零：

\[A\boldsymbol x=\boldsymbol0.\]

用第一节得到的阶梯形，它等价于

\[\begin{cases}
x_1+2x_2-x_3+x_4=0,\\
3x_3-3x_4=0.
\end{cases}\]

主元在第1、3列，所以第2、4个未知量是自由未知量。令它们分别为
s、t，从最后一式向上回代：

\[x_2=s,\quad x_4=t
\quad\Longrightarrow\quad
x_3=t,\quad x_1=-2s.\]

全部解便是

\[\boldsymbol x=
\begin{pmatrix}-2s\\s\\t\\t\end{pmatrix}
=s\underbrace{\begin{pmatrix}-2\\1\\0\\0\end{pmatrix}}_{\boldsymbol\xi_1}
+t\underbrace{\begin{pmatrix}0\\0\\1\\1\end{pmatrix}}_{\boldsymbol\xi_2}.\]

自由分量分别取
(1,0)、(0,1)，就得到了这两个基础解向量。在它们的零线性组合中比较第2、4个分量，两个系数只能都为零，所以线性无关；上式又表示了全部解，因此它们组成一个基础解系。

\Needspace{7\baselineskip}

\subsection{非齐次方程组的一个特解与全部解}\label{ux975eux9f50ux6b21ux65b9ux7a0bux7ec4ux7684ux4e00ux4e2aux7279ux89e3ux4e0eux5168ux90e8ux89e3}

保留 A，只把右端换成下面的常数向量：

\[A\boldsymbol x=\boldsymbol b,
\qquad\boldsymbol b=\begin{pmatrix}1\\5\\4\end{pmatrix}.\]

增广矩阵的常数列必须跟着行变换一起变。沿用前面的消元，再将第二行除以
3、第一行加上第二行，得到

\[\left(\begin{array}{rrrr|r}
1&2&-1&1&1\\2&4&1&-1&5\\1&2&2&-2&4
\end{array}\right)
\longrightarrow
\left(\begin{array}{rrrr|r}
1&2&-1&1&1\\0&0&3&-3&3\\0&0&0&0&0
\end{array}\right)
\longrightarrow
\left(\begin{array}{rrrr|r}
1&2&0&0&2\\0&0&1&-1&1\\0&0&0&0&0
\end{array}\right).\]

它对应两个有效方程：

\[x_1+2x_2=2,\qquad x_3-x_4=1.\]

令自由未知量 x2、x4
都为零，就得到一个特解；再把前面齐次方程组的任意解加上去，得到全部解：

\[\boxed{\boldsymbol x=
\underbrace{\begin{pmatrix}2\\0\\1\\0\end{pmatrix}}_{\text{一个特解}}
+s\begin{pmatrix}-2\\1\\0\\0\end{pmatrix}
+t\begin{pmatrix}0\\0\\1\\1\end{pmatrix}.}\]

把非齐次方程组的常数项全换成零，得到的齐次方程组叫作导出组。这里的基础解系属于\textbf{导出组}，由两个四维解向量组成；第二节从
A 中选出的则是三维列向量，两者是不同的对象。

为什么能这样相加？两个非齐次解之差满足右端为零的方程；反过来，一个特解加上任意齐次解，右端仍然是原来的
b。因此要写完整通解，既要找到特解，也要写齐导出组的自由部分。

\textbf{短检查3} 在刚才的非齐次通解中取
s=1、t=0，写出一个具体解，再代回原方程组检查。

\Needspace{7\baselineskip}

\subsection{有解与解的个数}\label{ux6709ux89e3ux4e0eux89e3ux7684ux4e2aux6570}

设 A 有 n 列，也就是有 n 个未知量，增广矩阵记为右侧添加常数列的矩阵。

\[\overline A=(A\mid\boldsymbol b),\qquad r=\operatorname{rank}A.\]

{\def\LTcaptype{none} % do not increment counter
\begin{longtable}[]{@{}ll@{}}
\toprule\noalign{}
消元后发现什么 & 结论 \\
\midrule\noalign{}
\endhead
\bottomrule\noalign{}
\endlastfoot
系数全为零，常数项非零的一行 & 无解 \\
系数秩与增广秩相等，且等于 n & 唯一解 \\
系数秩与增广秩相等，且 r\textless n & 无穷多解，有 n-r 个自由未知量 \\
\end{longtable}
}

齐次方程组总有零解。当 r\textless n 时有非零解，基础解系恰含 n-r
个向量；当 r=n 时只有零解。这里 n 是未知量数，不能换成方程数。

求向量的线性表示，也可以转成这项计算。本例中

\[\boldsymbol b=2\boldsymbol\alpha_1+\boldsymbol\alpha_3.\]

所以常数列能由系数列表示，方程组有解；相应系数就是一个特解
(2,0,1,0)。判断一个向量能否由一组向量表出，可以把它放在增广列，按同样的方法比较秩。

\Needspace{7\baselineskip}

\section{4
含参数时先把临界值分开}\label{4-ux542bux53c2ux6570ux65f6ux5148ux628aux4e34ux754cux503cux5206ux5f00}

前面的主元都是确定的非零数。含参数时，能否作除法也成了题目的一部分。先找可能为零的量，再分类，最后在各类里消元或回代。

\Needspace{7\baselineskip}

\subsection{对角元含参数的三元方程组}\label{ux5bf9ux89d2ux5143ux542bux53c2ux6570ux7684ux4e09ux5143ux65b9ux7a0bux7ec4}

讨论下面的方程组何时有解，并求出全部解。参数 lambda 属于数域 P。

\[\begin{cases}
\lambda x_1+x_2+x_3=1,\\
x_1+\lambda x_2+x_3=\lambda,\\
x_1+x_2+\lambda x_3=\lambda^2.
\end{cases}\]

先算系数行列式，沿第一行展开并分解：

\[D=\begin{vmatrix}\lambda&1&1\\1&\lambda&1\\1&1&\lambda\end{vmatrix}
=\lambda^3-3\lambda+2=(\lambda-1)^2(\lambda+2).\]

因此，lambda=1 与 lambda=-2 要单独处理。

\textbf{第一类 lambda既不等于1也不等于负2}
此时行列式非零，方程组有唯一解。为减少重复计算，记三个未知量的和为
S。三式相加，得到

\[(\lambda+2)S=1+\lambda+\lambda^2,
\qquad S=\frac{1+\lambda+\lambda^2}{\lambda+2}.\]

第一式还可以写成下面这样，其余两式同理：

\[(\lambda-1)x_1+S=1,\quad
(\lambda-1)x_2+S=\lambda,\quad
(\lambda-1)x_3+S=\lambda^2.\]

在当前这一类参数下，可以除以 lambda-1。分别代入 S 并化简，得到

\[\boxed{(x_1,x_2,x_3)=
\left(-\frac{\lambda+1}{\lambda+2},\
\frac1{\lambda+2},\
\frac{(\lambda+1)^2}{\lambda+2}\right).}\]

\textbf{第二类 lambda等于1} 三式都变成同一个方程，全部解为

\[x_1+x_2+x_3=1,
\qquad (x_1,x_2,x_3)=(1-s-t,s,t).\]

\textbf{第三类 lambda等于负2} 三式相加，左端为零，右端却为 3，所以无解：

\[0=1-2+4=3.\]

虽然一般解约分以后在 lambda=1 时也能代入，但推导中曾经除过
lambda-1。代入那个式子至多得到一个具体解，不能据此宣布唯一解。

\textbf{短检查4} 当 lambda=1
时，写出两个不同的解，亲自确认它们都满足原方程组。

\Needspace{7\baselineskip}

\subsection{从矛盾行读出两个参数的条件}\label{ux4eceux77dbux76feux884cux8bfbux51faux4e24ux4e2aux53c2ux6570ux7684ux6761ux4ef6}

下面给出教材第23题消元后的增广矩阵。求 a、b 何时有解，并在有解时写通解。

\[\left(\begin{array}{rrrrr|r}
1&1&1&1&1&1\\
0&1&2&2&6&3\\
0&0&0&0&0&a\\
0&0&0&0&0&b-2
\end{array}\right).\]

末两行分别要求 0=a、0=b-2，所以有解当且仅当
a=0、b=2。此时主元在第1、2列，其余三个未知量自由。令
x3=s、x4=t、x5=u，回代得到

\[x_2=3-2s-2t-6u,\qquad x_1=-2+s+t+5u.\]

全部解为

\[\begin{aligned}
\boldsymbol x={}&(-2,3,0,0,0)^{\mathsf T}
+s(1,-2,1,0,0)^{\mathsf T}\\
&+t(1,-2,0,1,0)^{\mathsf T}
+u(5,-6,0,0,1)^{\mathsf T}.
\end{aligned}\]

若 a 或 b 不满足条件，不必继续找特解：矛盾行已经说明无解。

\Needspace{7\baselineskip}

\section{5
行列式先找结构再计算}\label{5-ux884cux5217ux5f0fux5148ux627eux7ed3ux6784ux518dux8ba1ux7b97}

求秩只关心主元个数；求行列式还要跟踪换行的符号和倍乘的系数。做题时先看能否造零、是否是熟悉的特殊形式，再决定怎样展开。

\Needspace{7\baselineskip}

\subsection{造零后只展开一次}\label{ux9020ux96f6ux540eux53eaux5c55ux5f00ux4e00ux6b21}

很多行列式的第一步，是让某一行或某一列只剩一个非零元素。先看下面这个 n
阶矩阵，其中 n 至少为 2：

\[a_{ii}=i,\qquad a_{ij}=2\ (i\ne j),\qquad n\ge2.\]

第二行全是 2。让其他每一行减去第二行，第二行保持不动：

\[R_i\leftarrow R_i-R_2\quad(i\ne2).\]

行列式不变，变换后为

\[D_n=\begin{vmatrix}
-1&0&0&\cdots&0\\
2&2&2&\cdots&2\\
0&0&1&\cdots&0\\
\vdots&\vdots&\vdots&\ddots&\vdots\\
0&0&0&\cdots&n-2
\end{vmatrix}.\]

第2列现在只有一个非零元素
2。沿这一列展开，位置符号为正；删去第2行、第2列后，剩下的是对角行列式：

\[D_n=2\cdot(-1)^{2+2}\cdot(-1)\cdot1\cdot2\cdots(n-2)
=\boxed{-2(n-2)!}.\]

\textbf{负号来自第一行留下的 -1，阶乘来自剩下的连续整数乘积。} 当 n=2
时，直接得到 -2，与约定 0!=1 一致。

\textbf{短检查5} n=4
时，这个行列式等于多少？写成数值，并指出负号从哪里来。

\Needspace{7\baselineskip}

\subsection{认出范德蒙德的连续幂}\label{ux8ba4ux51faux8303ux5fb7ux8499ux5fb7ux7684ux8fdeux7eedux5e42}

当一列的元素是同一个数的连续幂时，可以先辨认这一列对应哪个参数。标准形式是

\[\begin{vmatrix}
1&1&\cdots&1\\
a_1&a_2&\cdots&a_n\\
a_1^2&a_2^2&\cdots&a_n^2\\
\vdots&\vdots&&\vdots\\
a_1^{n-1}&a_2^{n-1}&\cdots&a_n^{n-1}
\end{vmatrix}
=\prod_{1\le i<j\le n}(a_j-a_i).\]

乘积符号的意思是：每一对参数只取一次，按照\textbf{后面的参数减前面的参数}相乘。三阶时要写全三个差：

\[V=\begin{vmatrix}1&1&1\\1&2&5\\1&4&25\end{vmatrix}
=(2-1)(5-1)(5-2)=12.\]

先核对每列是否真的是 1、某个数、这个数的平方。不能只看第一行都是 1
就套公式。转置不改变行列式，交换两列则变号。

\textbf{短检查6} 把刚才 V
的前两列交换，结果是多少？这一问不需要重新展开。

\Needspace{7\baselineskip}

\subsection{把代数余子式的和还原成行列式}\label{ux628aux4ee3ux6570ux4f59ux5b50ux5f0fux7684ux548cux8fd8ux539fux6210ux884cux5217ux5f0f}

遇到一串带代数余子式的和，先看系数是否来自同一行。这里用 Cij
表示代数余子式，用 Mij 表示不带位置符号的余子式：

\[C_{ij}=(-1)^{i+j}M_{ij}.\]

固定第 i 行，把这行的代数余子式分别乘上一组数再相加，就等于把原第 i
行换成这组数后的行列式：

\[\sum_{j=1}^{n}b_jC_{ij}
=\det\bigl(\text{把原矩阵第 }i\text{ 行换成 }(b_1,\ldots,b_n)\bigr).\]

因此，用原第 i
行作系数，得到原行列式；用另一条原行作系数，换行后出现两条相同行，结果为零。

继续使用上一小节列参数为 1、2、5 的范德蒙德行列式 V。原第一行是
(1,1,1)，第二行是 (1,2,5)，所以

\[C_{11}+C_{12}+C_{13}=12,
\qquad C_{11}+2C_{12}+5C_{13}=0.\]

这里不必分别算出三个代数余子式。若题目给的是
Mij，先把位置符号补齐，不能直接套这两条式子。

\Needspace{7\baselineskip}

\subsection{分块三角时分别计算对角方块}\label{ux5206ux5757ux4e09ux89d2ux65f6ux5206ux522bux8ba1ux7b97ux5bf9ux89d2ux65b9ux5757}

左下角或右上角整块为零时，可以把两个对角方块分别计算。要求这两个对角块各自是方阵，但不要求它们可逆：

\[\det\begin{pmatrix}B&C\\0&D\end{pmatrix}=\det B\det D.\]

例如 B 为下面的二阶矩阵，D 也是二阶矩阵，上右角 C
可以任意，则整个四阶行列式为 24：

\[B=\begin{pmatrix}1&2\\0&3\end{pmatrix},\quad
D=\begin{pmatrix}2&1\\0&4\end{pmatrix},\quad
\det B\det D=(1\cdot3)(2\cdot4)=24.\]

一般的四块矩阵，不能直接按二阶行列式那样把四块相乘相减。

\Needspace{7\baselineskip}

\section{6
核对六个短检查}\label{6-ux6838ux5bf9ux516dux4e2aux77edux68c0ux67e5}

现在回头看六个短检查，先补出刚才做记号的那一步。答案不同的时候，沿着对应的那次行变换或回代找差别。

\textbf{检查1 非零行与主元列} 有两个非零行，主元在第1、3列，因此秩为 2。

\textbf{检查2 换一个极大无关组}
可以选原第2、4列。它们分别是第1列的两倍、第3列的相反数，非零倍乘不会破坏这两个向量的无关性；两个向量又能表出全部原列，所以构成另一极大无关组。

\textbf{检查3 取一组具体自由分量} s=1、t=0 时，具体解是

\[\boldsymbol x=(0,1,1,0)^{\mathsf T}.\]

代回三行，左端分别为 1、5、4，与原常数列一致。若你得到的是
(0,1,0,0)，就是漏掉了特解中第三个分量的 1。

\textbf{检查4 特殊参数下的不同解} lambda=1
时，可取下面两个不同的解；它们的分量和都为 1，因而都满足原来的三条方程：

\[(1,0,0)^{\mathsf T},\qquad(0,1,0)^{\mathsf T}.\]

\textbf{检查5 阶乘与负号} n=4 时，结果为 -2乘2!，即 -4。展开位置 (2,2)
的符号是正，负号来自剩下对角上的 -1。

\textbf{检查6 换列以后} 原值为 12，只交换一次两列，变成 -12。

\Needspace{7\baselineskip}

\subsection{最后五分钟核这五件事}\label{ux6700ux540eux4e94ux5206ux949fux6838ux8fd9ux4e94ux4ef6ux4e8b}

\begin{enumerate}
\def\labelenumi{\arabic{enumi}.}
\tightlist
\item
  求行列式时，交换了几次行或列？有没有漏记符号或提出的公因子？
\item
  含参数的除法，分母为零的情形是否已经单独讨论？
\item
  极大无关组是否取自原矩阵？若按主元列选取，是否把列号对应回原矩阵？
\item
  基础解系是否有 n-r 个向量，每个向量是否有 n 个分量？
\item
  写出的特解是否代回得到原常数列？每个基础解向量代回是否得到零列？
\end{enumerate}

如果还有时间，只重做两处：从空白纸重新消元第一节的 A，再重新计算第五节的
n
阶行列式。能够自己完成这两处，便把今晚最重要的消元和造零方法各练了一遍。

\Needspace{7\baselineskip}

\section{7 选读补充计算}\label{7-ux9009ux8bfbux8865ux5145ux8ba1ux7b97}

下面用于补查现有笔记中的其他计算形式。先完成前六节，再看这一节。

\Needspace{7\baselineskip}

\subsection{三对角行列式可以递推降阶}\label{ux4e09ux5bf9ux89d2ux884cux5217ux5f0fux53efux4ee5ux9012ux63a8ux964dux9636}

设 n 阶矩阵只有主对角及紧邻的两条对角线上可能非零，主对角全为
d，上方紧邻对角全为 u，下方紧邻对角全为 l。它的行列式满足

\[D_0=1,\qquad D_1=d,\qquad
D_n=dD_{n-1}-ulD_{n-2}\quad(n\ge2).\]

先算前两步，便能检查自己是否漏了负号：

\[D_2=d^2-ul,\qquad D_3=d^3-2dul.\]

\Needspace{7\baselineskip}

\subsection{排列符号与三角形行列式}\label{ux6392ux5217ux7b26ux53f7ux4e0eux4e09ux89d2ux5f62ux884cux5217ux5f0f}

排列的逆序数，是前面比后面大的数对的个数；行列式排列项的符号为 -1
的逆序数次幂。交换任意两个位置会改变奇偶性；相邻交换才保证逆序数只改变
1。

上三角或下三角行列式等于主对角元素的乘积。若矩阵\textbf{只有反对角线上可能非零}，把列顺序全部倒过来，就能化成对角矩阵；额外符号为

\[(-1)^{n(n-1)/2}.\]

\Needspace{7\baselineskip}

\subsection{实三维向量的几何计算}\label{ux5b9eux4e09ux7ef4ux5411ux91cfux7684ux51e0ux4f55ux8ba1ux7b97}

在实三维空间，两个向量正交时点积为零；三个向量的混合积是按这三个向量为行排成的三阶行列式，平行六面体体积取它的绝对值：

\[\boldsymbol u\cdot\boldsymbol v=\sum_{i=1}^{3}u_iv_i,\qquad
V=\left|\det\begin{pmatrix}\boldsymbol u\\\boldsymbol v\\\boldsymbol w\end{pmatrix}\right|.\]

叉积可以用下面的二阶余子式计算，第二个分量前有一个负号：

\[\boldsymbol u\times\boldsymbol v=
(u_2v_3-u_3v_2,\ u_3v_1-u_1v_3,\ u_1v_2-u_2v_1).\]

\Needspace{7\baselineskip}

\subsection{判断数域时检查非零数的倒数}\label{ux5224ux65adux6570ux57dfux65f6ux68c0ux67e5ux975eux96f6ux6570ux7684ux5012ux6570}

数域是复数域中包含
0、1，并对加、减、乘及除以非零数封闭的集合。例如整数集不是数域，因为 2
的倒数不在整数集中。检查一个集合是否为数域时，最后别漏掉``非零元素的倒数仍在集合里''这一项。

\Needspace{7\baselineskip}

\subsection{需要回查某一步时}\label{ux9700ux8981ux56deux67e5ux67d0ux4e00ux6b65ux65f6}

行列式方法对应原笔记 1.3至1.6，秩与方程组对应
2.1至2.3。详细解释可从下面的入口继续查，不必重新读完整套笔记：

\begin{itemize}
\tightlist
\item
  \href{https://www.xy0v0.top/notes/higher-algebra/02-02-matrix-rank}{矩阵的秩}
\item
  \href{https://www.xy0v0.top/notes/higher-algebra/02-03-linear-systems}{线性方程组的解}
\item
  \href{https://www.xy0v0.top/notes/higher-algebra/01-06-cramers-rule-and-vandermonde-determinants}{范德蒙德行列式与克拉默法则}
\end{itemize}

\end{document}
